\section{\textcolor{MidnightBlue}{Zustandsanalyse des alten Systems als Vorlage zur Anforderungsanalyse für das neue System}}
\label{ch:IST-Zustand:sec:Zustandsanalyse}

Die bestehende Oberfläche zur Anzeige der Testpool-Ergebnisse ist über die Jahre gewachsen und nicht mehr die optimale Lösung. Zur Vorbereitung der Anforderungsanalyse für das neue System werden nachfolgend Merkmale aufgelistet, welche beibehalten werden sollen, sowie Schwächen und Probleme identifiziert, welche durch das neue System zu beheben sind.

% **************************************************************************
%           funktionale Merkmale
% **************************************************************************
Das System zeigt alle Ergebnisse des Testpools vollständig und tagesaktuell an. Dabei können die Ergebnisse nach genutzer Maschine und Versionsnummer sowie nach Art des Tests unterschieden werden. Das System ermöglicht es, Details der Testdurchführung anhand von Logeinträgen einzusehen.

% **************************************************************************
%           Schwächen / Probleme
% **************************************************************************
Die Menge der Daten ermöglicht es dem alten System nicht, eine Übersichtlichkeit herzustellen. Das bedeutet, dass der Überblick über viele einzelne Oberflächen behalten werden muss und das System es nicht eindeutig ermöglicht, zwischen den Anwendungsfällen zu wechseln. Weiterhin sind die verfügbaren Oberflächen nicht ausreichend dokumentiert.

%- Identifizierte Schwächen / Probleme\\
%z.B. Usability-Probleme, technische Einschränkungen, fehlende Features
%Seiten / Oberflächen nicht gesammelt verfügbar

% **************************************************************************
%           ISO-Tabellen
% **************************************************************************
Die Umsetzung der Usability dieses Systems wurde anhand der \glqq Tabelle zur Unterstützung der Anwendung [der ISO 9241-110]\grqq{} \cite{DIN_EN_ISO_9241-110} bewertet. % S. 36ff.
Darin wird erfasst, welche Empfehlungen für diesen Anwendungsfall umsetzbar sind und ob diese umgesetzt wurden. Nachfolgend ist eine Zusammenfassung der Tabellen in \autoref{tab:Zusammenfassung:ISO-Tabellen:alt_Zustand} anhand der Interaktionsprinzipien aufgelistet.

%ISO-Tabelle:
%- nach Prinzips aufgesplittet (mit je der autoref)
%- Statististische Auswertung (einfach zahlen: a von b sind umsetzbar \& x von y wurden umgesetzt)

\begin{table}[H]
    \centering
    \begin{tabular}{|l|c|c|c|}
        \hline
        Prinzip Nr. & \makecell{Anzahl\\Empfehlungen} & \makecell{davon\\umsetzbar} & \makecell{davon\\umgesetzt} \\ \hline
        1) Aufgabenangemessenheit              & 6  & 6 & 2 \\ \hline
        2) Selbstbeschreibungsfähigkeit        & 8  & 6 & 3 \\ \hline
        3) Erwartungskonformität               & 11 & 8 & 5 \\ \hline
        4) Erlernbarkeit                       & 6  & 6 & 1 \\ \hline
        5) Steuerbarkeit                       & 10 & 6 & 4 \\ \hline
        6) Robustheit gegen Benutzungsfehler   & 12 & 1 & 0 \\ \hline
        7) Benutzerbindung                     & 12 & 4 & 3 \\ \hline
        \hline
        Summe                                  & 65 & 37 & 18 \\ \hline
    \end{tabular}
    \caption{Zusammenfassung der Umsetzung der Empfehlungen aus \autoref{Add:ISO-Tabellen} anhand des alten Systems}
    \label{tab:Zusammenfassung:ISO-Tabellen:alt_Zustand}
\end{table}
